% !TEX program = xelatex
% Türkçe akademik makale iskeleti — XeLaTeX + polyglossia + biblatex/biber
% Derleme:  python3 scripts/build.py makale-tr.tex
\documentclass[12pt,a4paper]{article}

% --- Dil ve font (Türkçe: XeLaTeX yolu; inputenc/fontenc KULLANILMAZ)
\usepackage{fontspec}
\usepackage{polyglossia}
\setmainlanguage{turkish}
\setotherlanguage{english}
\setmainfont{DejaVu Serif}[Scale=0.92]
\usepackage{csquotes}

% --- Yerleşim
\usepackage[a4paper,margin=2.5cm]{geometry}
\usepackage{setspace}
\onehalfspacing

% --- İçerik paketleri
\usepackage{amsmath,amssymb}
\usepackage{graphicx}
\graphicspath{{sekiller/}}
\usepackage{booktabs}
\usepackage{siunitx}
\sisetup{output-decimal-marker={,}, group-separator={.}}
\usepackage{subcaption}
\usepackage{enumitem}

% --- Kaynakça
\usepackage[backend=biber,style=numeric,sorting=none,maxbibnames=99]{biblatex}
\addbibresource{kaynaklar.bib}
% biblatex'in Türkçe "içinde" dizesi cümle başında \MakeCapital ile "Içinde" olur
% (doğrusu "İçinde"). Doğrulanmış düzeltme:
\newcommand{\bbxTRicinde}{İçinde}
\DefineBibliographyStrings{turkish}{%
  in = {{\noexpand\protect\noexpand\bbxTRicinde}{i\c{c}inde}},
}

% --- Referanslar (hyperref en son, cleveref ondan sonra)
\usepackage[hidelinks]{hyperref}
\usepackage{cleveref}
\crefname{figure}{Şekil}{Şekiller}
\crefname{table}{Tablo}{Tablolar}
\crefname{equation}{Denklem}{Denklemler}
\crefname{section}{Bölüm}{Bölümler}

\title{Makale Başlığı: Alt Başlık Gerekiyorsa}
\author{Ad Soyad\\ \small Kurum Adı, Bölüm \\ \small \texttt{ad@kurum.edu.tr}}
\date{\today}

\begin{document}
\maketitle

\begin{abstract}
\noindent Türkçe özet: problem, yöntem, temel bulgu, katkı. Tek paragraf, 150--250 kelime.
Atıf içermez, kısaltma tanımlanmadan kullanılmaz.

\vspace{4pt}
\noindent\textbf{Anahtar Kelimeler:} birinci, ikinci, üçüncü, dördüncü
\end{abstract}

\begin{english}
\begin{abstract}
\noindent English abstract goes here. Note that the language switch also changes
hyphenation rules, so wrapping it in an \texttt{english} environment matters.

\vspace{4pt}
\noindent\textbf{Keywords:} first, second, third, fourth
\end{abstract}
\end{english}

\section{Giriş}
\label{sec:giris}
Problem, bağlam ve neden önemli olduğu. Literatürdeki boşluk. Bu çalışmanın katkısı
madde madde. Örnek atıf~\cite{ornek2024}.

\section{İlgili Çalışmalar}
\label{sec:ilgili}
Tematik gruplayarak yaz, kronolojik liste yapma.

\section{Yöntem}
\label{sec:yontem}
Denklem örneği:
\begin{equation}
  \label{eq:kayip}
  \mathcal{L} = \frac{1}{N}\sum_{i=1}^{N}\left(y_i - \hat{y}_i\right)^2
\end{equation}
\Cref{eq:kayip} ortalama karesel hatayı verir. Ondalık sayı: \num{3,14}.

\section{Bulgular}
\label{sec:bulgular}

\begin{table}[htbp]
  \centering
  \caption{Yöntemlerin karşılaştırması. Başlık tabloda üstte olur.}
  \label{tab:sonuc}
  \begin{tabular}{l S[table-format=2.1] S[table-format=2.1]}
    \toprule
    Yöntem & {Doğruluk (\%)} & {Süre (s)} \\
    \midrule
    Temel çizgi & 82,4 & 12,1 \\
    Önerilen    & 91,7 &  9,8 \\
    \bottomrule
  \end{tabular}
\end{table}

\Cref{tab:sonuc} önerilen yöntemin üstünlüğünü gösteriyor.

% Şekil örneği (sekiller/ klasörüne dosya konunca yorum kaldırılır)
% \begin{figure}[htbp]
%   \centering
%   \includegraphics[width=0.8\linewidth]{sonuc.pdf}
%   \caption{Şekil başlığı altta olur.}
%   \label{fig:sonuc}
% \end{figure}

\section{Tartışma}
\label{sec:tartisma}
Bulguların anlamı, sınırlılıklar, tehditler.

\section{Sonuç}
\label{sec:sonuc}
Katkının özeti ve gelecek çalışma.

\section*{Teşekkür}
Fon ve destek bilgisi. \textbf{Anonim gönderimde bu bölüm kaldırılır.}

\printbibliography[title={Kaynaklar}]

\end{document}
